\section{Related Work}\label{sec:related}
Unsupervised time‑series clustering is a well‑studied problem.
Classical approaches include k‑means on engineered features, dynamic time warping (DTW) based clustering~\cite{Keogh2005}, and more recent deep learning methods that learn latent representations via autoencoders or contrastive learning~\cite{Franceschi2019}.
These methods typically require choosing a distance metric or a representation model that is domain‑specific.

Graph‑based community detection has been applied to time series by constructing networks from correlation or mutual information~\cite{Lacasa2008}.
However, the resulting communities are often interpreted post‑hoc and lack a formal framework linking them to functional modules.

The APEIRON pipeline differs in that it is not a clustering algorithm per se, but a layered architecture where the co‑activation graph (Layer~2) emerges naturally from the co‑occurrence statistics of elementary observables, and the functional clusters (Layer~3) are validated through their causal role in the system's predictive performance.
This provides a principled, domain‑independent methodology for unsupervised structure discovery that is grounded in the categorical foundations of the framework~\cite{Beniers2026Categorical}.